\chapter{Cronograma de actividades}

El presente capítulo describe la planificación temporal propuesta para el desarrollo del
proyecto. El cronograma se elaboró considerando la complejidad de las diferentes etapas
de diseño, implementación, integración y validación del sistema, así como los tiempos
necesarios para la elaboración de la documentación y la incorporación de observaciones
realizadas por el tribunal evaluador.

La planificación contempla un periodo académico de dieciséis semanas para el desarrollo
del proyecto y una semana adicional destinada a la presentación y defensa final ante el
tribunal evaluador. Se establece como objetivo alcanzar una versión funcional y validada
del prototipo al finalizar la Semana 15.

\section{Planificación general}

El desarrollo del proyecto se divide en siete etapas principales:

\begin{enumerate}
    \item Investigación bibliográfica y definición de requerimientos.
    \item Diseño de la arquitectura del sistema.
    \item Adquisición e integración del hardware.
    \item Desarrollo de software y de la interfaz gráfica.
    \item Implementación de mecanismos de identificación y almacenamiento.
    \item Integración y validación del sistema.
    \item Elaboración y revisión de la documentación final.
\end{enumerate}

La Figura o Tabla de cronograma resume la distribución temporal de las actividades
propuestas para el desarrollo del proyecto.

\section{Cronograma propuesto}

\begin{figure}[H]
        \centering
        \includegraphics[width=1\textwidth]{pics/SP_21.png}
        \caption{Cronograma propuesto, fuente: elaboración propia}
        \label{fig:ambar}
\end{figure}

%\begin{table}[H]
%\centering
%\caption{Cronograma general del proyecto.}
%\label{tab:cronograma}
%\renewcommand{\arraystretch}{1.2}
%\resizebox{\textwidth}{!}{
%\begin{tabular}{lccccccccccccccccc}
%\toprule
%\textbf{Actividad}
%& \textbf{1}
%& \textbf{2}
%& \textbf{3}
%& \textbf{4}
%& \textbf{5}
%& \textbf{6}
%& \textbf{7}
%& \textbf{8}
%& \textbf{9}
%& \textbf{10}
%& \textbf{11}
%& \textbf{12}
%& \textbf{13}
%& \textbf{14}
%& \textbf{15}
%& \textbf{16}
%& \textbf{17} \\
%\midrule
%
%Investigación bibliográfica
%& X & X & X & & & & & & & & & & & & & &\\
%
%Definición de requerimientos
%& X & X & X & & & & & & & & & & & & & &\\
%
%Retroalimentación y ajustes del alcance
%& X & X & X & & & & & & & & & & & & & &\\
%
%Diseño de arquitectura del sistema
%& & X & X & X & & & & & & & & & & & & &\\
%
%Adquisición de componentes
%& X & X & X & X & & & & & & & & & & & & &\\
%
%Implementación de identificación por código de barras
%& & & & X & X & X & & & & & & & & & & &\\
%
%Implementación de monitoreo ambiental
%& & & & X & X & & & & & & & & & & & &\\
%
%Implementación de reconocimiento facial
%& & & & & X & X & X & X & & & & & & & & &\\
%
%Implementación de almacenamiento de registros
%& & & & & & X & X & X & X & & & & & & & &\\
%
%Desarrollo de interfaz gráfica
%& & & & & & & & X & X & X & X & & & & & &\\
%
%Desarrollo del sistema de reportes
%& & & & & & & & X & X & X & X & & & & & &\\
%
%Integración hardware y software
%& & & & & & & & & X & X & X & X & & & & &\\
%
%Pruebas funcionales
%& & & & & & & & & & & & X & X & & & &\\
%
%Validación del sistema
%& & & & & & & & & & & & & X & X & X & &\\
%
%Ajustes finales del prototipo
%& & & & & & & & & & & & X & X & X & X & &\\
%
%Documentación del proyecto
%& X & X & X & X & X & X & X & X & X & X & X & X & X & X & X & X &\\
%
%Defensa del proyecto
%& & & & & & & & & & & & & & & & & X\\
%
%\bottomrule
%\end{tabular}
%}
%\end{table}

\section{Descripción de las etapas}

Durante las primeras semanas se realizará la investigación bibliográfica, la definición de
requerimientos y la consolidación del alcance del proyecto mediante reuniones periódicas
con el profesor asesor. Esta etapa permitirá establecer la arquitectura definitiva del sistema
y realizar los ajustes necesarios antes de iniciar la implementación.

Posteriormente se llevará a cabo la integración de los componentes de hardware y el
desarrollo de los diferentes módulos de software, incluyendo la interfaz gráfica, la
identificación de usuarios, el monitoreo ambiental, el almacenamiento de registros y la
generación de reportes.

Una vez implementados los subsistemas principales, se realizará la integración completa
del sistema y las pruebas funcionales correspondientes. La validación formal del prototipo
se efectuará entre las semanas 12 y 15, con el objetivo de verificar el cumplimiento de los
requerimientos definidos en el proyecto.

Se establece la semana 15 como fecha objetivo para alcanzar una versión funcional y
validada del prototipo. La semana 16 se reserva para la corrección de errores,
 incorporación de observaciones finales y actualización de la documentación técnica.

